\clearpage
\section{Consulta rápida}
\setlength{\LTpre}{4pt}
\setlength{\LTpost}{6pt}

As tabelas agrupam os comandos pelo assunto. Cada linha traz um comando ou expressão e sua função no contexto dos exemplos do guia.

\needspace{9\baselineskip}
\subsection{Python}
\subsubsection*{Operadores}

\begin{longtable}{@{}>{\raggedright\arraybackslash}p{0.44\linewidth}@{\hspace{1em}}>{\raggedright\arraybackslash}p{\dimexpr0.56\linewidth-1em\relax}@{}}
\toprule
\textbf{Comando ou expressão} & \textbf{Uso} \\
\midrule
\endfirsthead
\toprule
\textbf{Comando ou expressão} & \textbf{Uso} \\
\midrule
\endhead
\midrule
\endfoot
\bottomrule
\endlastfoot
\codigo{x + y} & Somar valores numéricos. \\
\codigo{x - y} & Subtrair \(y\) de \(x\). \\
\codigo{x * y} & Multiplicar valores numéricos. \\
\codigo{x / y} & Dividir e obter um resultado de ponto flutuante. \\
\codigo{x // y} & Dividir e arredondar o quociente para baixo. \\
\codigo{x \% y} & Obter o resto da divisão. \\
\codigo{x ** y} & Elevar \(x\) à potência \(y\). \\
\codigo{x = valor} & Associar um valor ao nome \codigo{x}. \\
\codigo{x += y} & Atualizar um número somando \(y\) ao valor anterior. \\
\codigo{x == y} & Testar se os valores são iguais. \\
\codigo{x != y} & Testar se os valores são diferentes. \\
\codigo{x < y} & Testar se \(x\) é menor que \(y\). \\
\codigo{x <= y} & Testar se \(x\) é menor ou igual a \(y\). \\
\codigo{x > y} & Testar se \(x\) é maior que \(y\). \\
\codigo{x >= y} & Testar se \(x\) é maior ou igual a \(y\). \\
\codigo{a and b} & Exigir as duas condições booleanas. \\
\codigo{a or b} & Exigir pelo menos uma condição booleana. \\
\codigo{not a} & Inverter uma condição booleana. \\
\codigo{x in valores} & Testar se um valor pertence à coleção. \\
\codigo{x not in valores} & Testar se um valor está fora da coleção. \\
\codigo{x is None} & Identificar o valor usado para representar ausência. \\
\codigo{x is not None} & Testar se o valor é diferente de \codigo{None} por identidade. \\
\end{longtable}

\needspace{6\baselineskip}
\subsubsection*{Tipos e texto}

\begin{longtable}{@{}>{\raggedright\arraybackslash}p{0.44\linewidth}@{\hspace{1em}}>{\raggedright\arraybackslash}p{\dimexpr0.56\linewidth-1em\relax}@{}}
\toprule
\textbf{Comando ou expressão} & \textbf{Uso} \\
\midrule
\endfirsthead
\toprule
\textbf{Comando ou expressão} & \textbf{Uso} \\
\midrule
\endhead
\midrule
\endfoot
\bottomrule
\endlastfoot
\codigo{print(x)} & Exibir um valor na saída. \\
\codigo{type(x)} & Consultar o tipo de um objeto. \\
\codigo{int(x)} & Converter um valor para inteiro. \\
\codigo{float(x)} & Converter um valor para ponto flutuante. \\
\codigo{str(x)} & Converter um valor para texto. \\
\codigo{bool(x)} & Converter um valor para booleano. \\
\codigo{texto[i]} & Selecionar o caractere na posição \(i\). \\
\codigo{texto[inicio:fim]} & Extrair um trecho, excluindo o índice final. \\
\codigo{texto.strip()} & Remover espaços das extremidades. \\
\codigo{texto.lower()} & Retornar uma versão em letras minúsculas. \\
\codigo{texto.replace(a, b)} & Substituir cada ocorrência de \codigo{a} por \codigo{b}. \\
\codigo{texto.split(",")} & Separar o texto em partes usando a vírgula. \\
\codigo{" ".join(partes)} & Juntar textos separados por um espaço. \\
\codigo{f"\{nota:.2f\}"} & Apresentar a nota com duas casas decimais. \\
\codigo{f"\{proporcao:.0\%\}"} & Apresentar uma proporção como porcentagem inteira. \\
\end{longtable}

\needspace{6\baselineskip}
\subsubsection*{Coleções e repetição}

\begin{longtable}{@{}>{\raggedright\arraybackslash}p{0.44\linewidth}@{\hspace{1em}}>{\raggedright\arraybackslash}p{\dimexpr0.56\linewidth-1em\relax}@{}}
\toprule
\textbf{Comando ou expressão} & \textbf{Uso} \\
\midrule
\endfirsthead
\toprule
\textbf{Comando ou expressão} & \textbf{Uso} \\
\midrule
\endhead
\midrule
\endfoot
\bottomrule
\endlastfoot
\codigo{len(valores)} & Contar os itens de uma coleção. \\
\codigo{sum(valores)} & Somar os valores de uma coleção numérica. \\
\codigo{min(valores)} & Obter o menor valor. \\
\codigo{max(valores)} & Obter o maior valor. \\
\codigo{sorted(valores)} & Criar uma lista ordenada. \\
\codigo{lista.sort()} & Ordenar a lista existente. \\
\codigo{lista.append(x)} & Acrescentar um item ao final da lista. \\
\codigo{lista.extend(xs)} & Acrescentar os itens de outra coleção à lista. \\
\codigo{matriz[i][j]} & Ler a coluna \(j\) da linha \(i\) em uma lista de listas. \\
\codigo{(a, b)} & Criar uma tupla com dois elementos. \\
\codigo{set(valores)} & Criar um conjunto de valores distintos. \\
\codigo{dicionario.get(chave, padrao)} & Ler uma chave e retornar o padrão se ela estiver ausente. \\
\codigo{dicionario.items()} & Percorrer pares de chave e valor. \\
\codigo{range(inicio, fim, passo)} & Gerar inteiros até antes do limite final. \\
\codigo{enumerate(valores)} & Percorrer cada item junto com sua posição. \\
\codigo{zip(a, b)} & Agrupar elementos correspondentes até acabar a sequência mais curta. \\
\end{longtable}

\needspace{6\baselineskip}
\subsection{Estatística e probabilidade}

\begin{longtable}{@{}>{\raggedright\arraybackslash}p{0.44\linewidth}@{\hspace{1em}}>{\raggedright\arraybackslash}p{\dimexpr0.56\linewidth-1em\relax}@{}}
\toprule
\textbf{Comando ou expressão} & \textbf{Uso} \\
\midrule
\endfirsthead
\toprule
\textbf{Comando ou expressão} & \textbf{Uso} \\
\midrule
\endhead
\midrule
\endfoot
\bottomrule
\endlastfoot
\codigo{np.unique(x, return\_counts=True)} & Obter valores distintos e suas contagens. \\
\codigo{rotulos[contagens == maior]} & Identificar as modas quando \codigo{maior} é a maior contagem. \\
\codigo{np.count\_nonzero(mascara)} & Contar as posições verdadeiras de uma máscara. \\
\codigo{round(x, 2)} & Arredondar um número para duas casas decimais. \\
\codigo{binom.pmf(k, n, p)} & Calcular a probabilidade de exatamente \(k\) sucessos em \(n\) tentativas binomiais. \\
\codigo{poisson.pmf(k, mu)} & Calcular a probabilidade de exatamente \(k\) ocorrências para a média informada. \\
\end{longtable}

\needspace{9\baselineskip}
\subsection{NumPy}
\subsubsection*{Criação, formato e seleção}

\begin{longtable}{@{}>{\raggedright\arraybackslash}p{0.44\linewidth}@{\hspace{1em}}>{\raggedright\arraybackslash}p{\dimexpr0.56\linewidth-1em\relax}@{}}
\toprule
\textbf{Comando ou expressão} & \textbf{Uso} \\
\midrule
\endfirsthead
\toprule
\textbf{Comando ou expressão} & \textbf{Uso} \\
\midrule
\endhead
\midrule
\endfoot
\bottomrule
\endlastfoot
\codigo{np.array(lista)} & Criar um array a partir de uma sequência. \\
\codigo{np.zeros((l, c))} & Criar uma matriz de zeros com \(l\) linhas e \(c\) colunas. \\
\codigo{np.ones(n)} & Criar um vetor de uns com \(n\) elementos. \\
\codigo{np.arange(a, b, p)} & Gerar valores com passo \(p\), excluindo \(b\) para passos inteiros. \\
\codigo{np.linspace(a, b, n)} & Gerar \(n\) pontos igualmente espaçados. \\
\codigo{arr.shape} & Consultar o tamanho de cada dimensão. \\
\codigo{arr.ndim} & Consultar o número de dimensões. \\
\codigo{arr.dtype} & Consultar o tipo dos elementos. \\
\codigo{arr.reshape(l, c)} & Reorganizar o array com \(l\) linhas e \(c\) colunas. \\
\codigo{arr[i, j]} & Ler a entrada da linha \(i\) e coluna \(j\). \\
\codigo{arr[:, 1]} & Selecionar a segunda coluna. \\
\codigo{arr[arr >= 7]} & Selecionar elementos de acordo com uma máscara booleana. \\
\codigo{a.T} & Obter a matriz transposta. \\
\codigo{a * b} & Multiplicar elementos em posições correspondentes de arrays compatíveis. \\
\codigo{a @ b} & Calcular o produto matricial. \\
\end{longtable}

\needspace{6\baselineskip}
\subsubsection*{Operações e estatísticas}

\begin{longtable}{@{}>{\raggedright\arraybackslash}p{0.44\linewidth}@{\hspace{1em}}>{\raggedright\arraybackslash}p{\dimexpr0.56\linewidth-1em\relax}@{}}
\toprule
\textbf{Comando ou expressão} & \textbf{Uso} \\
\midrule
\endfirsthead
\toprule
\textbf{Comando ou expressão} & \textbf{Uso} \\
\midrule
\endhead
\midrule
\endfoot
\bottomrule
\endlastfoot
\codigo{arr.sum(axis=0)} & Somar os valores de cada coluna. \\
\codigo{arr.sum(axis=1)} & Somar os valores de cada linha. \\
\codigo{np.min(arr)} & Obter o menor elemento. \\
\codigo{np.max(arr)} & Obter o maior elemento. \\
\codigo{np.mean(arr)} & Calcular a média. \\
\codigo{np.nanmean(arr)} & Calcular a média desconsiderando \codigo{NaN}. \\
\codigo{np.median(arr)} & Calcular a mediana. \\
\codigo{np.var(arr, ddof=1)} & Calcular a variância amostral com divisor \(n-1\). \\
\codigo{np.std(arr, ddof=1)} & Calcular o desvio padrão amostral com divisor \(n-1\). \\
\codigo{np.percentile(arr, q)} & Calcular o percentil \(q\). \\
\codigo{np.corrcoef(x, y)} & Calcular a matriz de correlação de Pearson. \\
\end{longtable}

\needspace{9\baselineskip}
\subsection{pandas}
\subsubsection*{Leitura e inspeção}

\begin{longtable}{@{}>{\raggedright\arraybackslash}p{0.44\linewidth}@{\hspace{1em}}>{\raggedright\arraybackslash}p{\dimexpr0.56\linewidth-1em\relax}@{}}
\toprule
\textbf{Comando ou expressão} & \textbf{Uso} \\
\midrule
\endfirsthead
\toprule
\textbf{Comando ou expressão} & \textbf{Uso} \\
\midrule
\endhead
\midrule
\endfoot
\bottomrule
\endlastfoot
\codigo{pd.read\_csv(caminho)} & Ler um arquivo CSV. \\
\codigo{pd.Series(valores)} & Criar uma série com valores e índice. \\
\codigo{pd.DataFrame(registros)} & Criar uma tabela a partir de registros. \\
\codigo{df.head(3)} & Exibir as três primeiras linhas. \\
\codigo{df.tail(2)} & Exibir as duas últimas linhas. \\
\codigo{df.info()} & Exibir tipos e contagens de valores presentes por coluna. \\
\codigo{df.shape} & Consultar número de linhas e colunas. \\
\codigo{df.columns} & Consultar os nomes das colunas. \\
\codigo{df.describe()} & Gerar um resumo estatístico das colunas numéricas. \\
\end{longtable}

\needspace{6\baselineskip}
\subsubsection*{Seleção, limpeza e transformação}

\begin{longtable}{@{}>{\raggedright\arraybackslash}p{0.44\linewidth}@{\hspace{1em}}>{\raggedright\arraybackslash}p{\dimexpr0.56\linewidth-1em\relax}@{}}
\toprule
\textbf{Comando ou expressão} & \textbf{Uso} \\
\midrule
\endfirsthead
\toprule
\textbf{Comando ou expressão} & \textbf{Uso} \\
\midrule
\endhead
\midrule
\endfoot
\bottomrule
\endlastfoot
\codigo{df["nota"]} & Selecionar uma coluna como série. \\
\codigo{df[["curso", "nota"]]} & Selecionar duas colunas como tabela. \\
\codigo{df.iloc[:3]} & Selecionar as três primeiras linhas pela posição. \\
\codigo{df.loc[0, "nota"]} & Selecionar a célula pelo rótulo da linha e da coluna. \\
\codigo{df.loc[mascara]} & Selecionar as linhas marcadas como verdadeiras. \\
\codigo{serie.isin(valores)} & Criar uma máscara de pertencimento. \\
\codigo{df.sort\_values("nota")} & Ordenar as linhas pela nota. \\
\codigo{df.isna()} & Identificar células com valores ausentes. \\
\codigo{mascara.sum()} & Contar valores verdadeiros por coluna de uma máscara booleana. \\
\codigo{df.dropna(subset=["nota"])} & Remover linhas com nota ausente. \\
\codigo{serie.fillna(valor)} & Preencher valores ausentes de uma série. \\
\codigo{pd.to\_numeric(serie, errors="coerce")} & Converter valores e sinalizar entradas inválidas como ausentes. \\
\codigo{serie.apply(funcao)} & Aplicar uma função a cada valor da série. \\
\codigo{df.copy()} & Criar uma cópia da tabela para análise separada. \\
\end{longtable}

\needspace{6\baselineskip}
\subsubsection*{Resumo, junção e exportação}

\begin{longtable}{@{}>{\raggedright\arraybackslash}p{0.44\linewidth}@{\hspace{1em}}>{\raggedright\arraybackslash}p{\dimexpr0.56\linewidth-1em\relax}@{}}
\toprule
\textbf{Comando ou expressão} & \textbf{Uso} \\
\midrule
\endfirsthead
\toprule
\textbf{Comando ou expressão} & \textbf{Uso} \\
\midrule
\endhead
\midrule
\endfoot
\bottomrule
\endlastfoot
\codigo{serie.value\_counts()} & Contar ocorrências de cada valor. \\
\codigo{df.groupby("curso")} & Agrupar as linhas de acordo com o curso. \\
\codigo{grupo.agg(medidas)} & Calcular as medidas indicadas para cada grupo. \\
\codigo{grupo.count()} & Contar valores presentes por grupo. \\
\codigo{grupo.size()} & Contar todas as linhas por grupo. \\
\codigo{serie.quantile(0.5)} & Calcular o quantil de ordem \(0{,}5\), a mediana. \\
\codigo{serie.corr(outra)} & Calcular a correlação de Pearson entre duas séries alinhadas pelo índice. \\
\codigo{df.merge(outro, on="id")} & Combinar tabelas pela chave \codigo{id}. \\
\codigo{df.reset\_index()} & Levar o índice antigo para uma coluna e criar um índice numérico. \\
\codigo{df.to\_csv("saida.csv", index=False)} & Exportar a tabela para CSV sem gravar o índice. \\
\end{longtable}

\needspace{9\baselineskip}
\subsection{Matplotlib}
\subsubsection*{Criação e leitura dos gráficos}

\begin{longtable}{@{}>{\raggedright\arraybackslash}p{0.44\linewidth}@{\hspace{1em}}>{\raggedright\arraybackslash}p{\dimexpr0.56\linewidth-1em\relax}@{}}
\toprule
\textbf{Comando ou expressão} & \textbf{Uso} \\
\midrule
\endfirsthead
\toprule
\textbf{Comando ou expressão} & \textbf{Uso} \\
\midrule
\endhead
\midrule
\endfoot
\bottomrule
\endlastfoot
\codigo{plt.subplots()} & Criar a figura e os eixos. \\
\codigo{ax.plot(x, y)} & Conectar os pontos na ordem fornecida. \\
\codigo{ax.axhline(y)} & Inserir uma linha horizontal de referência. \\
\codigo{ax.bar(categorias, alturas)} & Comparar valores por categoria. \\
\codigo{ax.bar\_label(barras)} & Exibir os valores junto às barras. \\
\codigo{ax.hist(valores, bins=limites)} & Contar observações em intervalos. \\
\codigo{ax.boxplot(valores)} & Resumir a distribuição com mediana, quartis e bigodes. \\
\codigo{ax.scatter(x, y)} & Mostrar a relação entre duas variáveis numéricas. \\
\codigo{ax.set(title=..., xlabel=..., ylabel=...)} & Definir título e rótulos dos eixos. \\
\codigo{ax.set\_xticks(posicoes, rotulos)} & Definir as posições e os rótulos no eixo horizontal. \\
\codigo{ax.grid(alpha=0.2)} & Exibir linhas de grade com transparência. \\
\codigo{label="Série"} & Nomear um elemento quando ele é desenhado. \\
\codigo{ax.legend(loc="upper left")} & Exibir a legenda no canto superior esquerdo. \\
\codigo{plt.show()} & Exibir a figura. \\
\codigo{fig.savefig("grafico.png", dpi=200)} & Salvar a figura em um arquivo. \\
\end{longtable}

\subsection{Onde continuar}

A documentação do Python em português reúne tutoriais e a referência da linguagem \cite{python}. As documentações de NumPy, pandas, Matplotlib e SciPy descrevem parâmetros e apresentam exemplos de uso \cite{numpy,pandas,matplotlib,scipy}. A documentação do Jupyter aborda a execução e a organização de notebooks \cite{jupyter}.
